% Generated draft: requires mathematical and source review.
\documentclass[11pt]{article}
\usepackage[margin=25mm]{geometry}
\usepackage{fontspec}
\setmainfont{Latin Modern Roman}
\usepackage{amsmath,amssymb,mathtools,unicode-math}
\setmathfont{Latin Modern Math}
\usepackage{enumitem,needspace,xurl,hyperref}
\setlength{\parindent}{0pt}
\setlength{\parskip}{5pt}
\setlength{\emergencystretch}{3em}
\newcommand{\N}{\mathbb{N}}
\newcommand{\Z}{\mathbb{Z}}
\newcommand{\Q}{\mathbb{Q}}
\newcommand{\R}{\mathbb{R}}
\newcommand{\C}{\mathbb{C}}
\newcommand{\F}{\mathbb{F}}
\newcommand{\A}{\mathbb{A}}
\newcommand{\PP}{\mathbb{P}}
\newcommand{\E}{\mathbb{E}}
\newcommand{\eps}{\varepsilon}
\newcommand{\dd}{\,\mathrm{d}}
\begin{document}
\section*{UnsolvedMath ID: 30004879}
\textbf{Source problem code:} catalogue entry
\subsection{Definitions and mathematical statement}
Let \(\mathfrak g\) be an affine Lie algebra in the scope intended by the source, and let \(U_q(\mathfrak g)\) denote its quantum enveloping algebra. The supplied record does not specify the ground field, restrictions on \(q\), the allowed affine types, or the precise evaluation-module construction. Let \(V(z)\) be a finite-dimensional evaluation module in that construction, with underlying module \(V\) and spectral parameter \(z\).

Write \(\operatorname{Aut}^{\mathrm{inv}}_{\mathrm{II}}(\mathfrak g)\) for the source's class of involutive automorphisms of the second kind. In particular, its elements satisfy \(\zeta^2=\mathrm{id}_{\mathfrak g}\). The exact definition of this class and the construction of the associated quantum automorphism \(\zeta_q\) are not supplied. Whenever \(\zeta_q\) is defined, the pullback \(\zeta_q^*V\) of a representation \(\rho:U_q(\mathfrak g)\to\operatorname{End}(V)\) has action
\[
 x\cdot_{\zeta_q^*V}v=\rho(\zeta_q(x))v.
\]

\textbf{Conjecture 3, as extracted from the supplied report.} For every evaluation module \(V(z)\) in the intended class, there exists
\[
 \zeta\in\operatorname{Aut}^{\mathrm{inv}}_{\mathrm{II}}(\mathfrak g)
 \quad\text{such that}\quad
 \zeta_q^*V\cong V.
\]
Here the displayed isomorphism interprets the source notation \(\zeta_q^*V=V\) as equivalence of modules. Under this interpretation, the conclusion means that there is an invertible linear map \(T:V\to V\) satisfying
\[
 T\rho(\zeta_q(x))=\rho(x)T
 \qquad\text{for every }x\in U_q(\mathfrak g).
\]
The automorphism \(\zeta\) may depend on \(V(z)\). Whether the source requires literal equality of actions instead of an isomorphism remains unresolved. The clean statement's word ``invariant'' is also unresolved: the original passage uses notation for involutions and discusses involutions of the second kind.

\textbf{Qualifications and motivation supplied by the source.} A pair of involutions of the second kind yields a solution of the universal reflection equation, labeled (2) in the report, in a completion of \(U_q(\mathfrak g)^{\otimes 2}\). Its evaluation on integrable category-\(\mathcal O\) modules is described in terms of a type-\(B\), or cylindrical, ribbon Artin braid groupoid. For finite-dimensional \(\mathfrak g\) and \(\zeta=\mathrm{id}\), the passage says this recovers the construction labeled [BK19]. In affine type, evaluation on finite-dimensional evaluation modules produces matrix-valued formal Laurent series. Under an unspecified ``mild assumption'' on \(\zeta\), the source asserts
\[
 \zeta_q^*(V(z))=(\zeta_q^*V)(z^{-1}),
\]
leading to the generalized parameter-dependent reflection equation labeled [Ch92]. The conjecture is intended to recover the original untwisted or transpose-twisted parameter-dependent reflection equation, with boundary quantum-integrability references [Ch84, Sk88]. The record does not state whether the unspecified mild assumption is part of the conjecture's required choice of \(\zeta\).

The supplied passage reports verification for \(U_q(\widehat{\mathfrak{sl}}_2)\). The supplied dated curation labels the problem partially solved and says that existence for arbitrary evaluation modules and affine types remains conjectural. These are supplied source assessments, not independently verified mathematical conclusions.
\subsection{Short English statement}
The conjecture asks whether every finite-dimensional quantum affine evaluation module admits an involutive automorphism of the second kind whose quantum pullback preserves the underlying module up to isomorphism. The exact evaluation-module class, quantum lift, parameter assumptions, and meaning of preservation require clarification. The supplied source reports the affine \(\mathfrak{sl}_2\) case as verified.
\subsection{Sources}
\begin{itemize}
\item Local UnsolvedMath registry: \texttt{data/unsolvedmath/problems.json}, permanent ID \(30004879\), source problem code \texttt{catalogue entry}; supplied complete source record. The record supplies no individual proposer.
\item Supplied authoritative report citation: \emph{Mini-Workshop: Three Facets of R-Matrices (2022)}, Oberwolfach Reports, pp.\ 2791--2825, DOI \texttt{10.4171/owr/2021/51}. \url{https://doi.org/10.4171/owr/2021/51}. The extracted passage labels the target Conjecture 3. The registry gives proposed year 2021; the supplied citation uses 2022.
\item Supplied literature source labeled ``Appel--Vlaar'': \url{https://arxiv.org/abs/2007.09218}. No title or precise supporting result is supplied in the record.
\item Literature assessment checked at \texttt{2026-08-22}: supplied dated curation, not independently verified. It labels the problem partially solved, reports verification of the reflection-equation conjecture for \(U_q(\widehat{\mathfrak{sl}}_2)\), and describes existence for arbitrary evaluation modules and affine types as conjectural. The background explicitly requests independent verification of time-sensitive status and cited claims.
\item The source passage also uses reference labels [BK19], [Ch92], [Ch84], and [Sk88]. Their bibliographic entries are absent from the supplied record; no identities or links are inferred here.
\item The registry labels the statement \texttt{corrected\_verified} and says it was checked against the official report. This is supplied verification metadata; this draft does not independently verify that claim or the cited sources.
\end{itemize}
\end{document}
